% Part: incompleteness
% Chapter: representability-in-q
% Section: c-representable

\documentclass[../../../include/open-logic-section]{subfiles}

\begin{document}

\olfileid[tr]{inc}{req}{cre}
\olsection{$C$ İçindeki Fonksiyonlar $\Th{Q}$'da Temsil Edilebilir}

$C$ içindeki her fonksiyonun $\Th{Q}$'da temsil edilebilir olduğunu
göstermeliyiz. Sonunda, $C$ içindeki her $k$-li
$f(x_0,\dots,x_{k-1})$ fonksiyonuna onu temsil eden bir
$!A_f(x_0,\dots,x_{k-1},y)$ formülünü nasıl atayacağımızı göstermemiz gerekir.

\begin{lem}
$n$ ve $m$ doğal sayılarsa ve $n \neq m$ ise
$Q \Proves \num n \neq \num m$.
\end{lem}

\begin{proof}
Her $m$ için $n \neq m$ ise $Q \Proves \eq/[\num n][\num m]$ olduğunu
$n$ üzerinde tümevarımla gösterin.

Temel durumda $n = 0$. $m$, $0$'a eşit değilse bir $k$ doğal sayısı için
$m = k + 1$ olur. $\lforall[x][\eq/[0][x']]$ diyen bir aksiyomumuz vardır.
Bir niceleyici aksiyomuyla $x$ yerine $\num k$ koyarak
$\eq/[0][\num k']$ sonucunu çıkarabiliriz. Fakat $\num k'$, yalnızca
$\num m$'dir.

Tümevarım adımında iddianın $n$ için doğru olduğunu varsayıp $n+1$'i ele
alabiliriz. $m$ herhangi bir doğal sayı olsun. İki olasılık vardır: ya $m=0$
ya da bir $k$ için $m=k+1$. İlk durum yukarıdaki gibi çözülür. İkinci durumda
$n+1 \neq k+1$ olduğunu varsayın. O zaman $n \neq k$. $n$ için tümevarım
varsayımına göre $\Th{Q} \Proves \eq/[\num n][\num k]$. Elimizde
$\lforall[x][\lforall[y][\eq[x'][y'] \lif \eq[x][y]]]$ diyen bir aksiyom
vardır. Bir niceleyici aksiyomuyla
$\eq[\num n'][\num k'] \lif \eq[\num n][\num k]$ elde ederiz. Önerme
mantığını kullanarak $\Th{Q}$ içinde
$\eq/[\num n][\num k] \lif \eq/[\num n'][\num k']$ çıkarımını yaparız.
Modus ponens ile $\eq/[\num n'][\num k']$ elde edilir; $\num k'$,
$\num m$ olduğundan istediğimiz de budur.
\end{proof}

Önermenin çok fazla şey söylemediğine dikkat edin: Özünde $\Th{Q}$'nun farklı
sayımların farklı nesneleri gösterdiğini ispatlayabildiğini söyler. Örneğin
$\Th{Q}$, $0'' \neq 0'''$ olduğunu ispatlar. Fakat bunun genel olarak
geçerli olduğunu göstermek biraz özen gerektirir. Tümevarım kullanıyor olsak
da bunun $\Th{Q}$'nun \emph{dışında} yapılan bir tümevarım olduğuna da dikkat
edin.

Sıfırı, ardılı, toplamayı, çarpmayı, eşitliğin karakteristik fonksiyonunu ve
izdüşümleri temsil edebileceğiz. Her durumda uygun temsil eden fonksiyon
son derece doğrudandır; örneğin sıfır
\[
y = 0,
\]
formülüyle, ardıl
\[
x_0' = y,
\]
formülüyle ve toplama
\[
x_0 + x_1 = y
\]
formülüyle temsil edilir. Yapılması gereken, $\Th{Q}$'nun ilgili ifadeleri
ispatlayabildiğini göstermektir. Örneğin toplamanın yukarıdaki formülle temsil
edildiğini söylemek, her $m$ ve $n$ doğal sayı çifti için $\Th{Q}$'nun
\[
\eq[\num n + \num m][\num {n+m}]
\]
ve
\[
\lforall[y][(\eq[(\num n + \num m)][y] \lif \eq[y][\num {n+m}])]
\]
ifadelerini ispatladığını göstermeyi gerektirir.

Peki ya bileşke? $h$ şu biçimde tanımlanmış olsun:
\[
h(x_0,\dots,x_{l-1}) = f(g_0(x_0,\dots,x_{l-1}), \dots,
g_{k-1}(x_0,\dots,x_{l-1})).
\]
$f,g_0,\dots,g_{k-1}$ fonksiyonlarını sırasıyla temsil eden
$!A_f, !A_{g_0}, \dots, !A_{g_{k-1}}$ formüllerini zaten bulduğumuzu
varsayın. O zaman $h$'yi temsil eden bir $!A_h$ formülünü,
$!A_h(x_0,\dots,x_{l-1},y)$ aşağıdaki formül olacak biçimde tanımlayabiliriz:
\begin{multline*}
  \lexists[z_0\dots][\lexists[z_{k-1}][(!A_{g_0}(x_0,\dots,x_{l-1},z_0) \land
  \dots \land !A_{g_{k-1}}(x_0,\dots,x_{l-1},z_{k-1}) \land]]\\
  !A_f(z_0,\dots,z_{k-1},y)).
\end{multline*}

Son olarak sınırsız aramayı ele alalım. $g(x,\vec z)$ düzenli ve
$\Th{Q}$'da, örneğin $!A_g(x,\vec z,y)$ formülüyle temsil edilebilir olsun.
$f$, $f(\vec z) = \mu x \; g(x,\vec z)$ ile tanımlansın. $f$'yi temsil eden
bir $!A_f(\vec z,y)$ formülü bulmak istiyoruz. Doğal bir seçim şudur:
\[
!A_f(\vec z,y) \ident !A_g(y,\vec z,0) \land \lforall[w][(w < y
\lif \lnot !A_g(w,\vec z,0))].
\]

Bunun işe yaradığı, örneğin aşağıdaki ek önermeler kullanılarak gösterilebilir.

\begin{lem}
  Her $x$ değişkeni ve her $n$ doğal sayısı için $\Th{Q}$,
  $\eq[(x' + \num n)][(x + \num n)']$ ifadesini ispatlar.
\end{lem}

Bunun, $\Th{Q}$'nun
$\lforall[x][\lforall[y][(x' + y) = (x +y)']]$ ifadesini ispatladığını
söylemekten daha zayıf olduğunu (ikincisinin yanlış olduğunu) bir kez daha
belirtmek gerekir.

\begin{proof}
İspat her zamanki gibi $n$ üzerinde tümevarımladır. $n=0$ temel durumunda
$\Th{Q}$'nun $\eq[(x'+0)][(x+0)']$ ifadesini ispatladığını göstermeliyiz.
Fakat elimizde şunlar vardır:
\begin{eqnarray*}
& & \eq[(x' + 0)][x'] \quad \text{4. aksiyomdan} \\
& & \eq[(x + 0)][x] \quad \text{4. aksiyomdan} \\
& & \eq[(x+0)'][x'] \quad \text{2. satırdan} \\
& & \eq[(x' + 0)][(x + 0)'] \quad \text{1. ve 3. satırlardan}
\end{eqnarray*}
Tümevarım adımında $\Th{Q}$ içinde
$\eq[(x' + \num n)][(x + \num n)']$ ifadesini türettiğimizi
varsayabiliriz. $\num{n+1}$, $\num n'$ olduğundan $\Th{Q}$'nun
$\eq[(x' + \num n')][(x + \num n')']$ ifadesini ispatladığını göstermeliyiz.
Aşağıdaki eşitlikler zinciri $\Th{Q}$ içinde türetilebilir:
\begin{eqnarray*}
(x' + \num n') & \eq & (x' + \num n)' \quad \text{5. aksiyom} \\
& \eq & (x + \num n')' \quad \text{tümevarım varsayımından}
\end{eqnarray*}
\end{proof}

\begin{lem}
\begin{enumerate}
\item $\Th{Q}$, $\lnot (x < \num 0)$ ifadesini ispatlar.
\item Her $n$ doğal sayısı için $\Th{Q}$ şunu ispatlar:
\[
x < \num {n+1} \lif (\eq[x][\num 0] \lor \dots \lor \eq[x][\num n]).
\]
\end{enumerate}
\end{lem}

\begin{proof}
1. maddeyi ve 2. maddenin bir kısmını gayriresmî olarak (yani biçimsel
!!{derivation}{}in nasıl kurulacağına ilişkin yalnızca ipuçları vererek)
yapalım.

1. madde için $<$ tanımı gereği $\Th{Q}$ içinde
$\lnot \lexists[y][\eq[(y'+x)][0]]$ ifadesini ispatlamamız gerekir. Bu,
birinci dereceden mantığın aksiyomları ve kuralları kullanılarak
$\lforall[y][\eq/[(y' + x)][0]]$ ifadesine denktir. Fikir şudur:
$\eq[(y'+x)][0]$ olsun. $x$, $0$ ise $\eq[(y' + 0)][0]$ elde ederiz. Fakat
$\Th{Q}$'nun 4. aksiyomuna göre $\eq[(y' + 0)][y']$, 2. aksiyomuna göre ise
$\eq/[y'][0]$ vardır; bu bir çelişkidir. Böylece
$\lforall[y][\eq/[(y' +x)][0]]$. $x$, $0$ değilse 3. aksiyoma göre
$\eq[x][z']$ olacak biçimde bir $z$ vardır. O zaman
$\eq[(y' + z')][0]$. 5. aksiyom uyarınca $\eq[(y' + z)'][0]$ elde edilir ve
bu yine 2. aksiyomla çelişir.

2. madde için $n$ üzerinde tümevarım kullanın. $n=0$ olan temel durumu ele
alalım. Bu durumda $x < \num 1 \lif \eq[x][\num 0]$ olduğunu göstermeliyiz.
$x < \num 1$ olsun. O zaman $<$ tanımlayan aksiyoma göre
$\lexists[y][\eq[(y'+x)][0']]$. $y$'nin bu özelliği taşıdığını varsayın;
dolayısıyla $\eq[y'+x][0']$.

$\eq[x][0]$ olduğunu göstermeliyiz. 3. aksiyoma göre $x$, $0$ değilse bir
$z$ için $z'$ye eşittir. O zaman $\eq[(y' + z')][0']$ elde ederiz.
$\Th{Q}$'nun 5. aksiyomuna göre $\eq[(y'+z)'][0']$, 1. aksiyomuna göre ise
$\eq[(y' + z)][0]$. Fakat bu tanım gereği $z < 0$ demektir ve 1. maddeyle
çelişir.
\end{proof}

\begin{explain}
Hesaplanabilir fonksiyonlar kümesinin, $\Th{Q}$'da temsil edilebilir
fonksiyonlar kümesi olarak nitelenebileceğini gösterdik. Gerçekte ispat daha
geneldir. Temsil edilebilirlik tanımından, $\Th{Q}$'yu genişleten (ya da
$\Th{Q}$'nun yorumlanabildiği) her kuramın hesaplanabilir fonksiyonları temsil
edebildiğini görmek zor değildir. Tersine, ispat kavramının hesaplanabilir
olduğu her !!{derivation} sisteminde temsil edilebilir her fonksiyon
hesaplanabilirdir. Dolayısıyla hesaplanabilir fonksiyonlar kümesi, örneğin
Peano aritmetiğinde, hatta Zermelo--Fraenkel küme kuramında temsil edilen
fonksiyonlar kümesi olarak nitelenebilir. Gödel'in belirttiği gibi bu biraz
şaşırtıcıdır. İspatlanabilirlik söz konusu olduğunda soruların ele alınan
kurama son derece duyarlı olduğunu göreceğiz; kabaca aksiyomlar ne kadar
güçlüyse o kadar çok şey ispatlanabilir. Buna karşılık geniş bir aksiyomatik
kuram yelpazesinde temsil edilebilir fonksiyonlar tam olarak hesaplanabilir
olanlardır.
\end{explain}

\end{document}
